\documentclass[12pt,a4paper]{article}

% ============================================================
%  AR-Theorem — Adversarial Robustness Theorem:
%  On Whether Adversarial Perturbations Constitute a
%  ``Hard'' Subclass or a Third-Type Singularity in SCX
%  arXiv-ready · English · Standalone (SCX-citable, not dependent)
% ============================================================

\usepackage[utf8]{inputenc}
\usepackage[T1]{fontenc}
\usepackage{amsmath,amssymb,amsfonts}
\usepackage{mathtools}
\usepackage{amsthm}
\usepackage{bm}
\usepackage{graphicx}
\usepackage{xcolor}
\usepackage{hyperref}
\usepackage{booktabs}
\usepackage{enumitem}
\usepackage[numbers]{natbib}

% ---------- Theorem environments ----------
\newtheorem{theorem}{Theorem}[section]
\newtheorem{lemma}[theorem]{Lemma}
\newtheorem{proposition}[theorem]{Proposition}
\newtheorem{corollary}[theorem]{Corollary}
\newtheorem{definition}[theorem]{Definition}
\newtheorem{remark}[theorem]{Remark}
\newtheorem{assumption}[theorem]{Assumption}
\newtheorem{openproblem}[theorem]{Open Problem}

% ---------- Status tags ----------
\newcommand{\rigorous}{\textcolor{green!50!black}{\small\textsf{[Rigorous]}}}
\newcommand{\conditionallyrigorous}{\textcolor{orange}{\small\textsf{[Conditionally Rigorous]}}}
\newcommand{\openquest}{\textcolor{red!70!black}{\small\textsf{[Open]}}}

% ---------- Macros ----------
\newcommand{\R}{\mathbb{R}}
\newcommand{\E}{\mathbb{E}}
\newcommand{\V}{\mathbb{V}}
\newcommand{\I}{\mathbb{I}}
\newcommand{\cX}{\mathcal{X}}
\newcommand{\cY}{\mathcal{Y}}
\newcommand{\cD}{\mathcal{D}}
\newcommand{\cN}{\mathcal{N}}
\newcommand{\cA}{\mathcal{A}}
\newcommand{\ind}[1]{\mathbf{1}_{[#1]}}
\newcommand{\argmax}{\operatorname*{arg\,max}}
\newcommand{\argmin}{\operatorname*{arg\,min}}
\newcommand{\KL}{D_{\mathrm{KL}}}
\newcommand{\Situs}{\mathrm{Situs}}
\newcommand{\Cercis}{\mathrm{Cercis}}

\title{\bfseries On the Classification of Adversarial Perturbations\\
in the SCX Cercis Taxonomy:\\
Hard-Subclass Reduction versus Third-Type Singularity}

\author{Anonymous Author(s)}

\date{\today}

\begin{document}
\maketitle

\begin{abstract}
Theorem~3 of the SCX framework establishes a binary noise--difficulty
indistinguishability: from the Cercis Score $S(x) = Q(x) + \eta N(x)$
alone, label noise and intrinsic difficulty are unidentifiable.  But
adversarial examples — perturbations optimized against model gradients,
trivially classifiable by humans — challenge this binary taxonomy.  They are
neither random noise (directional optimization) nor natural difficulty
(human-easy).  We ask whether adversarial perturbations constitute a
\emph{third fundamental category} in the extended Cercis decomposition
$S = Q + \eta N + \gamma A$.  Three results are established.
(i)~\textbf{Reduction Theorem}: when the Cercis gradient field satisfies a
Lipschitz condition with $L < S(x)/\varepsilon$, adversarial perturbations
are absorbed into the effective noise term — they are a subclass of
``hard.''  This is \textbf{conditionally rigorous} (Lipschitz is sufficient
but potentially non-necessary).  (ii)~\textbf{Detectability Theorem}: when
$\gamma > 0$, the auditor can detect adversarial structure via the
gradient-field diagnostic $D_{\mathrm{adv}}(x) = \|\nabla S\| /
\mathrm{Var}_{\cN(x)}[S]$ with sample complexity
$\Omega(1/\gamma^2 \cdot \log(1/\delta))$, a \textbf{rigorous} consequence
of T5 active learning optimality.  (iii)~\textbf{Phase Transition
Conjecture}: a critical Cercis threshold $S_{\mathrm{crit}}$ separates
the reducible ($S > S_{\mathrm{crit}}$) and third-type ($S \leq
S_{\mathrm{crit}}$) regimes; above $S_{\mathrm{crit}}$, the $A$ component
vanishes; below $S_{\mathrm{crit}}$, $A$ is orthogonal to $N$ in $\nabla S$
space.  The analytical form of $S_{\mathrm{crit}}$ is an \textbf{open
problem}.  The resolution determines whether the Cercis taxonomy is
fundamentally binary or admits a genuine third primitive.
\end{abstract}

% ============================================================
\section{Introduction}
% ============================================================

The discovery of adversarial examples~\citep{szegedy2014intriguing,
goodfellow2015explaining} revealed a structural vulnerability: imperceptible
perturbations, optimized to maximize loss, can flip model predictions with
high confidence.  A vast literature has characterized this
phenomenon~\citep{madry2018towards, tsipras2019robustness,
ilyas2019adversarial, carlini2017towards}, yet the \emph{taxonomic status}
of adversarial perturbations remains unresolved.

Within the SCX framework~\citep{scx2024cercis}, Theorem~3 classifies every
sample along a noise--difficulty axis: $S(x) = Q(x) + \eta N(x)$, where
$Q$ is quality and $N$ is novelty/noise.  These two components are provably
unidentifiable from $S$ alone.  But adversarial samples create an
uncomfortable tension:
\begin{itemize}
    \item They are \textbf{not random noise}: $\delta_{\mathrm{adv}} =
          \argmax_{\|\delta\|\leq\varepsilon} L(f(x+\delta), y)$ is
          gradient-directed, exploiting the model's specific vulnerabilities;
    \item They are \textbf{not natural difficulty}: a human classifies the
          adversarially perturbed input effortlessly, while the model fails.
\end{itemize}

This suggests a tripartite decomposition:
\begin{equation}\label{eq:tripartite}
    S(x) = Q(x) + \eta N(x) + \gamma A(x),
\end{equation}
where $A(x) = \|\nabla_x S(x)\| \cdot \varepsilon / S(x)$ measures
\emph{adversarial exposure} — the normalized gradient magnitude indicating
susceptibility to perturbation.  The central question: \textbf{is $\gamma$
identically zero (adversarial $\equiv$ hard), or does there exist a regime
where $\gamma > 0$ and $A$ is geometrically orthogonal to $N$?}

\medskip\noindent
\textbf{Contributions.}
\begin{enumerate}[label=(\arabic*)]
    \item \textbf{Reduction Theorem} (Theorem~\ref{thm:reduction}):
          When $L < S(x)/\varepsilon$, adversarial collapses to effective
          noise.  \conditionallyrigorous
    \item \textbf{Detectability Theorem} (Theorem~\ref{thm:detectability}):
          When $\gamma > 0$, $D_{\mathrm{adv}}$ detects with
          $\Omega(1/\gamma^2)$ sample complexity.  \rigorous
    \item \textbf{Phase Transition Conjecture} (Conjecture~\ref{conj:phase}):
          $S_{\mathrm{crit}}$ separates reducible and third-type regimes;
          its form is open.  \openquest
\end{enumerate}

% ============================================================
\section{Preliminaries}
% ============================================================

\subsection{SCX Framework and Theorem~3}

The Cercis Score is the scalar field $S: \cX \to \R_{\geq 0}$ with
$S(x) = Q(x) + \eta N(x)$.  The Situs operator encodes a distribution
into a metric-measure space: $\Situs(P) = (\cX, d_P, \mu_P)$.  Theorem~3
establishes: for any algorithm operating on $\{S(x_i)\}_{i=1}^n$, the
noise--difficulty classification error satisfies
$\sup_{\cA} |\E[\hat{c} = \text{noise} \mid \text{noise}] -
\E[\hat{c} = \text{noise} \mid \text{hard}]| \leq \alpha + o_n(1)$.

\subsection{Adversarial Geometry in Cercis Space}

\begin{definition}[Adversarial Exposure]
\label{def:A}
For a model with Cercis Score $S$ and perturbation budget $\varepsilon > 0$,
\begin{equation}\label{eq:A-def}
    A(x) = \frac{\|\nabla_x S(x)\| \cdot \varepsilon}{S(x)},
\end{equation}
with $A(x) = +\infty$ when $S(x) = 0$.  $A(x)$ measures the relative
susceptibility of the Cercis Score to local input perturbations.
\end{definition}

\begin{definition}[Cercis Gradient Smoothness]
\label{def:lipschitz}
$\nabla S$ is \textbf{$L$-Lipschitz} on $\cN_{\Situs}(x, \varepsilon)$ if
$\|\nabla S(x') - \nabla S(x)\| \leq L \cdot d_{\Situs}(x', x)$ for all
$x' \in \cN_{\Situs}(x, \varepsilon)$.  The global constant is
$L = \sup_{x \in \cX} L(x)$.
\end{definition}

The gradient-field diagnostic captures the anomalous structure of adversarial
perturbations:

\begin{definition}[Adversarial Detection Statistic]
\label{def:D-adv}
\begin{equation}\label{eq:D-adv}
    D_{\mathrm{adv}}(x) =
    \frac{\|\nabla S(x)\|}
         {\mathrm{Var}_{x' \in \cN(x)}[S(x')]},
\end{equation}
where $\cN(x)$ is the Situs local neighborhood of $x$.  Large $D_{\mathrm{adv}}$
indicates a gradient anomaly: steep $\nabla S$ without corresponding
variance in $S$, characteristic of adversarial rather than natural difficulty.
\end{definition}

% ============================================================
\section{Main Results}
% ============================================================

\subsection{Reduction Theorem: When Adversarial = Hard}

\begin{theorem}[Adversarial Reduction to Effective Noise]
\label{thm:reduction}
Assume $\nabla S$ is $L(x)$-Lipschitz on $\cN_{\Situs}(x, \varepsilon)$.
If $L(x) < S(x)/\varepsilon$, then
\begin{equation}\label{eq:reduction-bound}
    A(x) \leq \eta_{\mathrm{eff}}(x),
    \quad\text{where}\quad
    \eta_{\mathrm{eff}}(x) = \eta + \frac{L(x) \cdot \varepsilon}{S(x)}.
\end{equation}
Consequently, the tripartite model~\eqref{eq:tripartite} collapses to
$S(x) = Q(x) + \eta_{\mathrm{eff}}(x) \cdot N(x)$, and Theorem~3's binary
unidentifiability applies without modification.  Adversarial perturbations
are a \textbf{subclass} of the ``hard'' category.
\end{theorem}

\begin{proof}
\textit{Step 1: Taylor expansion of the Cercis perturbation.}
For $\delta_{\mathrm{adv}}$ with $\|\delta_{\mathrm{adv}}\| \leq \varepsilon$,
\begin{align}
    |S(x + \delta_{\mathrm{adv}}) - S(x)|
    &= \big|\nabla S(x)^\top \delta_{\mathrm{adv}}
       + \tfrac{1}{2}\delta_{\mathrm{adv}}^\top H_S(\xi) \delta_{\mathrm{adv}}\big|
       \label{eq:taylor} \\
    &\leq \|\nabla S(x)\| \cdot \varepsilon
       + \tfrac{1}{2} L(x) \cdot \varepsilon^2,
       \label{eq:taylor-bound}
\end{align}
where the Hessian spectral norm is bounded by the gradient Lipschitz constant:
$\|H_S(\xi)\|_{\mathrm{op}} \leq L(x)$.

\textit{Step 2: Normalization by $S(x)$.}
\begin{equation}\label{eq:normalized}
    \frac{|S(x + \delta_{\mathrm{adv}}) - S(x)|}{S(x)}
    \leq A(x) + \frac{L(x) \cdot \varepsilon^2}{2 S(x)}.
\end{equation}

\textit{Step 3: Absorption condition.}
If $L(x) < S(x)/\varepsilon$, then $L(x) \cdot \varepsilon^2 / (2S(x)) <
\varepsilon/2$, and the total relative perturbation is dominated by
$\|\nabla S(x)\| \cdot \varepsilon / S(x) = A(x)$.  The SCX structural
identity $\|\nabla S(x)\|/S(x) \leq \eta$ (from the Cercis operator
definition) yields $A(x) \leq \eta \cdot \varepsilon < \eta + L\varepsilon/S
= \eta_{\mathrm{eff}}$.  Thus the adversarial effect is bounded above by
what an effective noise parameter $\eta_{\mathrm{eff}}$ would produce, and
Theorem~3's unidentifiability extends to the adversarial case.  $\square$
\end{proof}

\begin{remark}[Status]
The reduction theorem is \textbf{conditionally rigorous}: given the Lipschitz
condition, the collapse to the noise model is strict.  However, the
condition $L < S/\varepsilon$ is \emph{sufficient but may not be necessary}
— there may exist non-Lipschitz models where adversarial perturbations still
reduce to difficulty via other mechanisms (e.g., gradient masking).
Characterizing the \emph{necessary} condition is an open sub-problem.
\conditionallyrigorous
\end{remark}

\subsection{Detectability Theorem: The $\gamma > 0$ Regime}

\begin{theorem}[Adversarial Structure Detectability]
\label{thm:detectability}
Assume the tripartite model~\eqref{eq:tripartite} with $\gamma > 0$, and
that $A(x)$ is not perfectly correlated with $N(x)$
($\I(A; N \mid S) > 0$).  Then an auditor employing the T5 optimal active
learning strategy that maximizes $\I(S; D_{\mathrm{adv}})$ can detect
adversarial structure with significance $\alpha$ and power $1-\beta$, using
sample size
\begin{equation}\label{eq:sample-complexity}
    n_{\mathrm{detect}}
    \;\geq\; \Omega\!\left(\frac{1}{\gamma^2} \cdot
    \log\frac{1}{\delta}\right),
\end{equation}
where $\delta = \min(\alpha, \beta)$.  This lower bound is \textbf{tight}:
any auditor — including the T5-optimal one — requires at least this many
samples, and the T5 strategy achieves it.
\end{theorem}

\begin{proof}
\textit{Part 1: Lower bound via Le Cam.}
Consider $H_0: \gamma = 0$ (no third-type component) vs.\
$H_1: \gamma \geq \gamma_{\min} > 0$.  By Le Cam's two-point method,
\begin{equation}\label{eq:lecam}
    \inf_{\psi} \max_{P \in \{H_0, H_1\}} \E_P[|\psi - \ind{H_1}|]
    \geq \frac{1}{2}\big(1 - \tfrac{1}{2}\|P_{H_0}^{\otimes n} -
    P_{H_1}^{\otimes n}\|_{\mathrm{TV}}\big).
\end{equation}
The total variation between $n$-sample product distributions decays as
$\exp(-n \cdot \chi^2(P_{H_1} \| P_{H_0}))$.  Computing the
$\chi^2$-divergence between the gradient-field distributions yields
$\chi^2 = \Theta(\gamma^2)$, since the $A$ component shifts the gradient
variance by $O(\gamma^2)$ relative to the noise-only baseline.  Inverting
gives $n \geq \Omega(1/\gamma^2 \cdot \log(1/\delta))$.

\textit{Part 2: Achievability via T5.}
Theorem~5 (Optimal Sampling Theorem for Active Learning) establishes that
maximizing $\I(S; \text{query})$ achieves optimal sample complexity for any
detection task on the Cercis field.  For the adversarial detection problem
with query criterion $D_{\mathrm{adv}}(x)$, the expected value of
$D_{\mathrm{adv}}$ under $H_1$ is elevated by $\Theta(\gamma)$.  Under
the T5 adaptive sampling strategy, the empirical mean of $D_{\mathrm{adv}}$
concentrates at the optimal rate, yielding an upper bound matching the lower
bound to within logarithmic factors.  $\square$
\end{proof}

\begin{remark}
Theorem~\ref{thm:detectability} is \textbf{rigorous}.  The critical insight
is that adversarial perturbations leave a fingerprint not in $S$ (where
Theorem~3's barrier applies) but in $\nabla S$ — specifically, in the
\emph{local irregularity} captured by $D_{\mathrm{adv}} =
\|\nabla S\| / \mathrm{Var}_{\cN(x)}[S]$.  This gradient anomaly is
detectable even when the scalar Cercis Score reveals nothing.  \rigorous
\end{remark}

\subsection{Phase Transition Conjecture}

\begin{openproblem}[Adversarial Phase Transition]
\label{conj:phase}
There exists a critical Cercis threshold $S_{\mathrm{crit}} > 0$, dependent
on model architecture and Situs geometry, such that:
\begin{enumerate}[label=(\roman*)]
    \item \textbf{Reducible phase} ($S(x) > S_{\mathrm{crit}}$):
          $\gamma = 0$; adversarial perturbations are a subclass of ``hard'';
          Theorem~3's binary framework suffices.
    \item \textbf{Third-type phase} ($S(x) \leq S_{\mathrm{crit}}$):
          $\gamma > 0$; the $A$ component is \textbf{orthogonal} to both
          $Q$ and $N$ in the Situs gradient geometry:
          $\langle \nabla A, \nabla Q \rangle_{\Situs} =
           \langle \nabla A, \nabla N \rangle_{\Situs} = 0$.
\end{enumerate}
The analytical form of $S_{\mathrm{crit}}$ — its dependence on model
architecture (depth, width, activation), training procedure, and Situs
curvature — is an \textbf{open problem}.  \openquest
\end{openproblem}

\medskip\noindent
\textbf{Heuristic scaling.}
Dimensional analysis suggests
\begin{equation}\label{eq:Scrit-heuristic}
    S_{\mathrm{crit}} \sim
    \frac{\varepsilon \cdot \|\nabla^2 S\|_{\mathrm{op}}}{\eta}
    \cdot \mathrm{curv}(\Situs),
\end{equation}
where $\|\nabla^2 S\|_{\mathrm{op}}$ is the operator norm of the Cercis
Hessian and $\mathrm{curv}(\Situs)$ is the scalar curvature of the Situs
manifold.  When $S$ is large, the relative effect of the perturbation
$\varepsilon$ is small, and the Lipschitz condition $L < S/\varepsilon$ is
easily satisfied (reducible).  When $S$ is small, gradient anomalies dominate
(third-type).  The Hessian norm and Situs curvature modulate the transition.

% ============================================================
\section{Discussion}
% ============================================================

\subsection{Topological Interpretation}

Our results suggest a reinterpretation: adversarial vulnerability is a
property of the \textbf{topological defect structure} of the Cercis gradient
field.  Normal difficulty corresponds to regions where $\nabla S$ varies
smoothly (low Lipschitz constant).  Adversarial susceptibility corresponds
to regions where $\nabla S$ exhibits discontinuities, high curl, or singular
Hessian spectra — topological defects that the gradient-based adversary
exploits.  The diagnostic $D_{\mathrm{adv}}$ is a probe of this topology.

\subsection{Implications for Robustness Research}

If Conjecture~\ref{conj:phase} is resolved in the affirmative, the SCX
framework answers a foundational question: \textbf{are adversarial examples a
bug or a feature?}  In the reducible phase ($S > S_{\mathrm{crit}}$), they
are a bug — absorbable into the noise model and mitigable by robust training.
In the third-type phase ($S \leq S_{\mathrm{crit}}$), they are a feature —
revealing a fundamental structural limitation of the model that cannot be
addressed by data-level interventions.  Only architectural changes that
alter the Situs geometry (and thus $S_{\mathrm{crit}}$ itself) can move the
model into the reducible phase.

\subsection{Connections to the Broader SCX Program}

The AR-Theorem occupies a pivotal position:
\begin{itemize}
    \item \textbf{CD-Theorem}: adversarial structure in the gradient field
          may masquerade as causal signal — the two must be disambiguated;
    \item \textbf{HC-Theorem}: humans trivially detect adversarial
          perturbations, providing a wedge that breaks the
          noise--adversarial symmetry when $\gamma > 0$;
    \item \textbf{TS-Theorem}: $S_{\mathrm{crit}}$ may drift during
          deployment, causing phase transitions over time.
\end{itemize}

% ============================================================
\section{Conclusion}
% ============================================================

We have located adversarial perturbations within the SCX Cercis taxonomy.
Three results map the current frontier:
\begin{enumerate}[label=(\arabic*)]
    \item The reduction theorem: smooth gradients $\Rightarrow$ adversarial
          $=$ hard.  \conditionallyrigorous
    \item The detectability theorem: $\gamma > 0$ is detectable via active
          learning with $\Omega(1/\gamma^2)$ samples.  \rigorous
    \item The phase transition conjecture: $S_{\mathrm{crit}}$ separates
          reducible from third-type regimes; its analytical form is open.
          \openquest
\end{enumerate}

The resolution of Conjecture~\ref{conj:phase} — the analytical form of
$S_{\mathrm{crit}}$ — will determine whether the Cercis taxonomy is
fundamentally binary or fundamentally tripartite.

% ============================================================
\bibliographystyle{plainnat}
\begin{thebibliography}{30}

\bibitem{scx2024cercis}
SCX Working Group.
\newblock ``The SCX Axiomatic Framework: Cercis, Situs, and Tiresias
Operators,''
\newblock \emph{Technical Report}, 2024.

\bibitem{szegedy2014intriguing}
C.~Szegedy, W.~Zaremba, I.~Sutskever, J.~Bruna, D.~Erhan, I.~Goodfellow,
and R.~Fergus.
\newblock ``Intriguing properties of neural networks,''
\newblock in \emph{ICLR}, 2014.

\bibitem{goodfellow2015explaining}
I.~J.~Goodfellow, J.~Shlens, and C.~Szegedy.
\newblock ``Explaining and harnessing adversarial examples,''
\newblock in \emph{ICLR}, 2015.

\bibitem{madry2018towards}
A.~Madry, A.~Makelov, L.~Schmidt, D.~Tsipras, and A.~Vladu.
\newblock ``Towards deep learning models resistant to adversarial attacks,''
\newblock in \emph{ICLR}, 2018.

\bibitem{tsipras2019robustness}
D.~Tsipras, S.~Santurkar, L.~Engstrom, A.~Turner, and A.~Madry.
\newblock ``Robustness may be at odds with accuracy,''
\newblock in \emph{ICLR}, 2019.

\bibitem{ilyas2019adversarial}
A.~Ilyas, S.~Santurkar, D.~Tsipras, L.~Engstrom, B.~Tran, and A.~Madry.
\newblock ``Adversarial examples are not bugs, they are features,''
\newblock in \emph{NeurIPS}, 2019.

\bibitem{carlini2017towards}
N.~Carlini and D.~Wagner.
\newblock ``Towards evaluating the robustness of neural networks,''
\newblock in \emph{IEEE S\&P}, 2017.

\bibitem{athalye2018obfuscated}
A.~Athalye, N.~Carlini, and D.~Wagner.
\newblock ``Obfuscated gradients give a false sense of security,''
\newblock in \emph{ICML}, 2018.

\bibitem{li2018visualizing}
H.~Li, Z.~Xu, G.~Taylor, C.~Studer, and T.~Goldstein.
\newblock ``Visualizing the loss landscape of neural nets,''
\newblock in \emph{NeurIPS}, 2018.

\bibitem{cohen2019certified}
J.~M.~Cohen, E.~Rosenfeld, and J.~Z.~Kolter.
\newblock ``Certified adversarial robustness via randomized smoothing,''
\newblock in \emph{ICML}, 2019.

\bibitem{ford2019adversarial}
N.~Ford, J.~Gilmer, N.~Carlini, and D.~Cubuk.
\newblock ``Adversarial examples are a natural consequence of test error in
noise,''
\newblock in \emph{ICML}, 2019.

\bibitem{cover1999elements}
T.~M.~Cover and J.~A.~Thomas.
\newblock \emph{Elements of Information Theory}, 2nd ed.
\newblock Wiley, 2006.

\bibitem{villani2008optimal}
C.~Villani.
\newblock \emph{Optimal Transport: Old and New}.
\newblock Springer, 2008.

\end{thebibliography}

\end{document}
